\documentclass[11pt]{article}
\usepackage[letterpaper,margin=1in]{geometry}
\usepackage[T1]{fontenc}
\usepackage{lmodern,microtype,mathtools,amssymb,amsthm,booktabs}
\usepackage{xcolor,fancyhdr}
\usepackage[colorlinks=true,allcolors=blue!45!black]{hyperref}
\pagestyle{fancy}\fancyhf{}
\fancyhead[L]{\small Indicator switching for pair coverage}\fancyhead[R]{\small Research draft}
\fancyfoot[C]{\thepage}\setlength{\headheight}{14pt}
\newtheorem{theorem}{Theorem}[section]
\newtheorem{lemma}[theorem]{Lemma}
\newtheorem{prop}[theorem]{Proposition}
\newtheorem{cor}[theorem]{Corollary}
\newcommand{\E}{\mathbb E}
\newcommand{\Prb}{\mathbb P}
\newcommand{\Var}{\operatorname{Var}}
\newcommand{\Cov}{\operatorname{Cov}}
\newcommand{\Poi}{\operatorname{Poi}}
\newcommand{\ind}{\mathbf1}
\newcommand{\dK}{d_{\mathrm K}}
\newcommand{\Nlaw}{\mathcal N(0,1)}
\newcommand{\sd}{\operatorname{sd}}
\newcommand{\kk}{\kappa}
\numberwithin{equation}{section}
\setlength{\emergencystretch}{1em}
\title{Indicator switching for Poissonized pair coverage}
\author{Christopher D. Long}
\date{November 27, 2026}
\begin{document}
\maketitle
\begin{abstract}
For arbitrary overlapping Poisson pair batches and arbitrary integer thresholds, a monotone switch
of a single coverage indicator gives a variance identity and an aggregate quadratic displacement
bound $3W V$. This bound includes the fluctuations of every affected neighbor.
It decreases to $2WV$ on triangle-free graphs. A separate expansion gives
$|\kappa_3(Y)|\le4WV$ on the full class. We prove the threshold-uniform Kolmogorov bound
$CW/\sqrt V$ on arbitrary unions of stars, allowing heterogeneous component sizes,
variances and thresholds. The full pair-graph Kolmogorov problem remains open.
\end{abstract}

\section{Setting and statements}
Independent clocks count the occurrences of each singleton and each unordered pair by a fixed
time, with Poisson laws of finite total mean. The graph consists of the positive-mean pair clocks.
Let $N_i$ be the total of clocks touching vertex $i$, $I_i=\ind_{\{N_i\ge r_i\}}$ for arbitrary
$r_i\in\{1,2,\ldots\}$, $p_i=\E I_i$, $v_i=p_i(1-p_i)$, $X_i=h_i(I_i-p_i)$,
$Y=\sum_i h_i I_i$, and $V=\Var Y>0$. Weights are nonnegative and their total on each
positive-mean batch is at most $W$. Inactive species can be deleted. Let $v_0=\sum_i\E X_i^2$.
Increasing functions of independent clocks are positively associated, so $v_0\le V$.
The elementary product association inequality follows by conditioning and the one-coordinate
identity $2\Cov(f(X),g(X))=\E[(f(X)-f(X'))(g(X)-g(X'))]$.

\begin{theorem}\label{t:cost}
There is a monotone coupling of $Y_i^+,Y_i^-$ having the laws of $Y$ given
$I_i=1,0$, respectively. With $D_i=Y_i^+-Y_i^-$, the same coupling satisfies
\begin{align}
 \Cov(Y,f(Y))&=\sum_i h_i v_i\E[f(Y_i^+)-f(Y_i^-)],\qquad
 V=\sum_i h_i v_i\E D_i,\label{e:identity}\\
 \sum_i h_i v_i\E D_i^2&\le \tfrac32 W(V+v_0)\le3WV.\label{e:cost}
\end{align}
For triangle-free graphs the right side of \eqref{e:cost} improves to $W(3V+v_0)/2$.
Vertices with $v_i=0$ can be omitted. The covariance identity holds, in particular, for
bounded Lipschitz tests and for the identity function.
\end{theorem}
A hub of threshold $n$ attached to $n$ unit-rate first-discovery leaves can move order
$\sqrt n$ leaves in conditional expectation (see the preceding manuscript \cite{previous}).
Equation \eqref{e:cost} bounds the sum of the actual second moments in just this regime.
It gives the correct dimensional order: disjoint equal-weight singleton indicators have
$\sum h_i v_i\E D_i^2=WV$ when their weights are $W$.

\begin{theorem}\label{t:block}
Suppose the pair graph is a disjoint union of stars. Then
\[
 \dK\big((Y-\E Y)/\sqrt V,\Nlaw\big)\le\min(1,CW/\sqrt V)
\]
with an absolute constant, uniformly over thresholds, rates, and the number and sizes of
the stars. More generally, suppose that within each connected component, deleting a set
of species of total weight at most $H$ leaves independent indicators within that component.
Then the bound holds with $C\max(W,H)/\sqrt V$.
\end{theorem}
Deleting a separate hub in every component may require unbounded total weight over the
whole graph. Theorem~\ref{t:block} covers this case even when each star has a different law.

\section{A direct coupling and its quadratic cost}
Let $K,K'$ be independent Poisson variables of mean $\lambda_i=\E N_i$, independent of the
nonincident clocks and of a common iid sequence of marks. Each mark selects an incident clock
with probability proportional to its mean. In the first environment use the first $K$ marks,
and in the second use the first $K'$ marks. All nonincident clocks agree. Each environment
has the original law. Conditional on $K\ge r_i>K'$ these are the two required conditional
environments. Write $s_i=|\ind_{\{K\ge r_i\}}-\ind_{\{K'\ge r_i\}}|$ and $L_j$ for the
absolute change of $X_j$, for $j\sim i$. Changes of all the vertex indicators have the same
sign as $K-K'$. With $L=h_i s_i+\sum_{j\sim i}L_j$ we have
\begin{equation}\label{e:sym}
 v_i\E D_i^m=\tfrac12\E[s_i L^m],\qquad m=1,2.
\end{equation}
Splitting $\Cov(I_i,f(Y))$ into its two conditional means proves \eqref{e:identity}.

We use two elementary projections of $X_j$. For an edge $ij$ with count $Z_{ij}$ let
$g_{j;i}(x)=\E[X_j\mid Z_{ij}=x]$. For two distinct edges at $j$ put
\[
 w_{j;ik}(x,u)=\E[X_j\mid Z_{ij}=x,Z_{jk}=u]-g_{j;i}(x)-g_{j;k}(u).
\]
Each $w$ has mean zero when averaged in either argument. Independence makes all distinct
first-order and unordered second-order projections at $j$ pairwise orthogonal; their
inner products with $X_j$ are their squared norms. In particular,
\begin{equation}\label{e:bessel}
 \sum_i\|g_{j;i}\|_2^2+\sum_{\{i,k\}}\|w_{j;ik}\|_2^2\le\E X_j^2 .
\end{equation}
Singleton clocks and higher projections can be left out of this inequality.

\begin{proof}[Proof of Theorem~\ref{t:cost}]
In $s_iL^2$ first take the squared self term, the self--neighbor terms and the squared
neighbor terms. Since $L_j^2=h_j L_j$ and $\E[s_i L_j]/2=\Cov(I_i,X_j)$ by monotonicity,
their total after summation with weight $h_i/2$ is
\[
 \sum_i h_i\Var X_i+\sum_i\sum_{j\sim i}(2h_i+h_j)\Cov(X_i,X_j)
 \le W v_0+\tfrac32 W(V-v_0).
\]
Indeed each unordered edge has coefficient $3(h_i+h_j)\le3W$, and all nonedge
covariances vanish. For distinct neighbors let $dX_j$ be the signed difference of $X_j$ in the two
environments. Then $L_jL_k=dX_jdX_k\ge0$, so $s_i$ can be discarded when taking an upper
bound on their total.

Condition on the two totals and the mark list. If $j,k$ share no edge, the external
clocks in $X_j$ and $X_k$ are disjoint, and the conditional product expectation is
$dg_{j;i}\,dg_{k;i}$. Here $d$ takes the difference between the two values of the
$ij$ or $ik$ count. If $jk$ is an edge, condition also on its count. The only extra
term after integrating that count is
\[
 \E_{Z_{jk}}[(d w_{j;ik})(d w_{k;ij})].
\]
In fact the $g_{j;k}(Z_{jk})$ part cancels on differencing, and the $dw$ terms have
zero mean in $Z_{jk}$. No other external clock can touch both $j$ and $k$.

For the first-order terms the sum over distinct $j,k$ is at most
$\E(\sum_j dg_{j;i})^2$. In each one of the two environments the edge counts out of
$i$ are independent Poisson by splitting. Moreover, conditional on the common marks, the two evaluations are iid as
functions of $K,K'$. Thus the expectation of their squared difference is twice the
expected conditional variance, at most $2\sum_j\|g_{j;i}\|_2^2$.
Including $h_i/2$ gives the first-order budget $W\sum_{j,i}\|g_{j;i}\|_2^2$.

For each ordered triangle the additional expectation is bounded in absolute value by
$2\|w_{j;ik}\|_2\|w_{k;ij}\|_2$. Indeed, conditioning on the common marks and
the $jk$ clock again gives two iid evaluations for each squared difference.
Use Cauchy--Schwarz, $2ab\le a^2+b^2$, and $h_i\le W$. In the ordered triples
a second-order projection with specified center $j$ and unordered neighbors
$\{i,k\}$ occurs twice in the sum of $\|w_{j;ik}\|_2^2$. With the $1/2$ in
\eqref{e:sym} the second-order budget is at most
$2W\sum_j\sum_{\{i,k\}}\|w_{j;ik}\|_2^2$.
By \eqref{e:bessel} the two budgets together are at most $2Wv_0$, or $Wv_0$ in
a triangle-free graph. Adding the self and repeated terms proves the constants.
For countable species, restrict to finite species sets. The switches can be realized
by projecting the same marks (an edge leaving the set becomes a singleton). Each
retained conditional displacement increases along an exhaustion. Monotone convergence,
and domination of $Y$ by $W$ times the total Poisson clock count for the variance
limit, complete the proof of the cost estimates. The identity for bounded
$f$ also follows by summing $h_i\Cov(I_i,f(Y))$ directly since
$\sum_i h_i p_i=\E Y<\infty$. For $f(y)=y$, use monotone convergence
of the nonnegative weighted covariances and $L^2$ convergence of the partial sums.
\end{proof}

\paragraph{The mean displacement has a separate bound.}
For all real $c_i$ the product-space martingale differences obtained by revealing a
clock touch at most two indicators. Cauchy--Schwarz and orthogonality give
$\Var(\sum_i c_i (I_i-p_i)/\sqrt{v_i})\le2\sum_i c_i^2$.
Duality gives $\sum_i\Cov(I_i,Y)^2/v_i\le2V$. Thus
$\sum_i h_i v_i(\E D_i)^2\le2WV$ directly from the conditional means.
The second-order projections in Theorem~\ref{t:cost} control what Jensen's inequality
alone leaves open.

\section{The cubic cumulant on arbitrary pair graphs}
\begin{theorem}\label{t:cubic}
For all pair graphs and arbitrary thresholds, $|\kappa_3(Y)|\le4WV$.
\end{theorem}
\begin{proof}
Expand $\E(\sum_i X_i)^3$. The equal-index terms cost at most $Wv_0$.
The terms with exactly two equal indices cost at most $\tfrac32 W(V-v_0)$ because
$\E X_i^2X_j=h_i(1-2p_i)\Cov(X_i,X_j)$, each unordered adjacent pair has total weight $h_i+h_j\le W$, and the covariances are nonnegative.

For three distinct vertices integrate out every clock except edges among the
three. Each conditioned vertex factor is the sum of its available first-order
projections and, if both adjacent edges occur, its second-order projection.
A term in the product integrates to zero if an edge used at one endpoint has
no factor at the other. The remaining possibilities are a path of two used edges
and a triangle of three used edges. For a path with center $j$, the expectation is
$\E[X_jg_{i;j}g_{k;j}]$; replacing the projected center by $X_j$ here is exact.
For fixed $j$ and ordered distinct ends, the bracket below has mean zero. Subtract the midpoint of the range of $X_j$ to obtain the bound
\[
 \tfrac12 W\E\left|(\sum_i g_{i;j})^2-\sum_i g_{i;j}^2\right|
 \le W\sum_i\E g_{i;j}^2 .
\]
There are three choices of position for the center in the ordered expansion.
The total path contribution is at most $3W\sum_{j,i}\|g_{i;j}\|_2^2$.

Integration over the three triangle edges gives
$|\E w_i w_j w_k|\le\|w_i\|_2\|w_j\|_2\|w_k\|_2$.
This follows by Cauchy--Schwarz in the edge shared by the first two factors
and then in the other two edges. Bound the product by the mean of the three cubed
norms. Each norm is at most $\|X_i\|_2\le W$ at its own vertex.
Summing over ordered triples costs at most $2W\sum_j\sum_{\{i,k\}}\|w_{j;ik}\|_2^2$. By \eqref{e:bessel} paths and triangles together cost at most $3Wv_0$.
Combining gives $4Wv_0+\tfrac32W(V-v_0)\le4WV$. Countable exhaustion is valid in $L^3$.
\end{proof}
The fourth cumulant need not be $O(W^2V)$ even on a star \cite{previous}.
Theorem~\ref{t:cubic} identifies a genuine change between orders three and four.

\section{Composition of independent star components}\label{sec:blocks}
We prove Theorem~\ref{t:block} as a stability result for independent components.
Write their centered contributions as $Z_b=S_b+R_b$, with $S_b$ the sum of the
independent centered weighted indicators left after deletion, and $R_b$ the
centered weight of the deleted indicators. Put $a=\max(W,H)$, $u_b=\Var Z_b$,
$s_b^2=\Var S_b$, and $\sigma^2=\sum_b u_b$. We have $|R_b|\le a$,
$s_b^2,\E R_b^2\le u_b$ by association, and independence between components.
The proof only needs these conditions and centered Bernoulli summands of range at most $a$.
We use $C$ for numerical constants which may change between inequalities.

\paragraph{One component.}
Expansion of the fourth power of an independent centered sum gives
$\E S_b^4\le3s_b^4+a^2s_b^2$; also $\E R_b^4\le a^2u_b$. Hence
\begin{equation}\label{e:moments}
 \E Z_b^4\le C(u_b^2+a^2u_b),\qquad \E|Z_b|^3\le C(\sqrt{u_b}+a)u_b .
\end{equation}
Its Kolmogorov error against a centered Gaussian of variance $u_b>0$ is at most
$Ca/\sqrt{u_b}$. Indeed if $\sqrt{u_b}\ge2a$, then
$|s_b-\sqrt{u_b}|\le a$, and independent Berry--Esseen for $S_b$ costs at most
$Ca/s_b$. The cdf sandwich from $|R_b|\le a$ then costs at most
$Ca/s_b$ for normal shifts and $C|\log(s_b/\sqrt{u_b})|$ for rescaling.
For $\sqrt{u_b}<2a$ use one.

There is also a smooth bound for $G_b\sim N(0,u_b)$:
\begin{equation}\label{e:smooth}
 |\E g(Z_b)-\E g(G_b)|\le Ca u_b\|g'''\|_\infty .
\end{equation}
Normalize $\|g'''\|_\infty=1$ and subtract the Taylor polynomial through degree two
at the origin; $Z_b,G_b$ have matching first two moments. Then $|g'(x)|\le x^2/2$.
Replacing $S_b+R_b$ by $S_b$ costs at most $Ca\E(S_b^2+R_b^2)\le Ca u_b$.
Replace each Bernoulli summand of $S_b$ by a centered Gaussian of the same variance.
Taylor expansion conditional on the other terms costs at most $Ca$ times the
summand variance. Finally couple $s_bG$ and $\sqrt{u_b}G$ by the same standard
normal $G$. Their scale difference is at most $a$, giving another $Ca u_b$ by the
derivative bound. This proves \eqref{e:smooth}.

\paragraph{Interval concentration for a subcollection.}
For $Q(X,l)=\sup_x\Prb(x\le X\le x+l)$ and a nonconstant sum $Z_J=\sum_{b\in J}Z_b$,
\begin{equation}\label{e:concentration}
 Q(Z_J,l)\le C(l+a)/\sqrt{\sum_{b\in J}u_b},\qquad l\ge0.
\end{equation}
Here we use the Miroshnikov--Rogozin inequality \cite{MR} in the form of
\cite[(1.2)--(1.3)]{BC}: for independent $Z_b$ and $L\ge a$ its upper bound is
$CL[\sum_{b\in J}m_b Q(Z_b,a)^{-2}]^{-1/2}$, where
$m_b=\E\min(|Z_b-Z_b'|,a)^2$ and $Z_b'$ is an independent copy.
To see that each contribution dominates $c u_b$, write $D=Z_b-Z_b'$.
Cauchy--Schwarz gives
$(2u_b)^2\le m_b\E(D^2+D^4/a^2)$ by factoring $D^2$ with $\min(|D|,a)$.
Equation \eqref{e:moments} yields $m_b\ge c u_b a^2/(a^2+u_b)$.
The one-component cdf estimate and the Gaussian density bound give
$Q(Z_b,a)\le\min(1,Ca/\sqrt{u_b})$, proving the assertion. Take $L=\max(l,a)$.

\begin{proof}[Proof of Theorem~\ref{t:block}]
Rescale so that $\sigma=1$ and henceforth let $a$ denote the rescaled bound.
It suffices to consider $a\le1$. If some $u_b>1/2$, replace that component by $G_b$.
By convolution contraction this costs $Ca$ in Kolmogorov distance. Its Gaussian
cdf smooths the test for all remaining components; the third derivative with respect
to their sum is bounded by $C u_b^{-3/2}\le C$. Successive applications of
\eqref{e:smooth} replace them by matching independent Gaussians at total cost
$Ca\sum_{j\ne b}u_j\le Ca$.

Otherwise put $U=\sum_b Z_b$ and $O_b=U-Z_b$.
The independent outside sum $O_b$ has variance at least $1/2$ for every $b$.
Thus \eqref{e:concentration} gives $Q(O_b,l),Q(U,l)\le C(l+a)$, the former also
conditional on the entire component $b$. For the normal Stein solution $f=f_x$ use
$|f|\le2$, $|f'|\le1$, and
\begin{equation}\label{e:steinvariation}
 |f'(y)-f'(u)|\le (2+|u|)|y-u|+
 \ind_{\{x\text{ lies between }u\text{ and }y\}} .
\end{equation}
This follows from $f'(u)=uf(u)+\ind_{\{u\le x\}}-\Phi(x)$, with the left
derivative at $x$.

Decompose the contribution of component $b$ to $\E U f(U)$ into
\begin{align*}
 &\E[S_b f(O_b+S_b)],\\
 &\E[S_b(f(U)-f(U-R_b))],\\
 &\E[R_b(f(U)-f(O_b))].
\end{align*}
Their sum is exact because $R_b$ is centered and independent of $O_b$.
For the first term write $S_b=\sum_j L_j$, $L_j=h_j(B_j-p_j)$ with independent
Bernoulli $B_j$. The two-point identity
\[
 \E[L_j f(O_b+S_b)]
 =\Var(L_j)\E f'(O_b+S_b-L_j+h_j(T-p_j))
\]
uses an independent uniform $T$ on $[0,1]$. The argument differs from $U$ by at
most $2a$. Equation \eqref{e:steinvariation}, $\E|U|\le1$ and $Q(U,4a)\le Ca$
therefore replace the first term by $s_b^2\E f'(U)$ with error $Ca s_b^2$.

In the second term integrate $f'$ along the segment of length $|R_b|$. Replacing
that derivative by $f'(U)$ gives $\E[S_bR_b f'(U)]$ with error at most $Ca u_b$.
For the smooth part of \eqref{e:steinvariation}, use $\E|S_bR_b|\le u_b$,
$\E|O_b|\le1$, and $\E|S_bR_bZ_b|\le a u_b$.
For the jump condition on block $b$: the interval for $O_b$ has length at most $a$
and probability at most $Ca$.

In the third term the segment has length $|Z_b|$. Its approximation by
$\E[R_b Z_b f'(U)]$ also costs at most $Ca u_b$. Indeed the smooth part costs at
most $\E[|R_b|Z_b^2(2+|O_b|+|Z_b|)]\le Ca u_b$ by independence and
\eqref{e:moments}. Conditional on this block the jump probability is at most
$C(|Z_b|+a)$, so its contribution is bounded by
$C\E[|R_bZ_b|(|Z_b|+a)]\le Ca u_b$.

The resulting multiplier of $f'(U)$ is
$M=\sum_b(s_b^2+S_bR_b+R_bZ_b)$. Its mean is $\sum_b u_b=1$.
Independence of the blocks and $|R_b|\le a$ give
$\Var M\le a^2\sum_b\E(S_b+Z_b)^2\le4a^2$.
Since $|f'|\le1$, this error and the three summed errors give
$|\E[f'(U)-Uf(U)]|\le Ca$. Taking the supremum over $x$ and undoing the
normalization proves the claim. Finite vertex exhaustion preserves the decomposition
by intersecting each deleted set with the retained species, using the original
independent component groups. Domination by the total Poisson count gives the
countable extension.
\end{proof}

\section{The general coupling remainder}
Choose the upper or lower endpoint of the $i$th coupling by an independent coin
with probabilities $p_i,1-p_i$. Denote the chosen environment by $\widehat\eta_i$
and its count by $\widehat Y_i$. Before weighting by $D_i$, both have the original law.
For any $f$ with Lipschitz derivative, \eqref{e:identity} gives
\begin{align}
 \Cov(Y,f(Y))-V\E f'(Y)
 &=\sum_i h_i v_i\Cov(D_i,f'(\widehat Y_i))+\mathcal R_f,\notag\\
 |\mathcal R_f|&\le3WV\operatorname{Lip}(f').\label{e:gap}
\end{align}
Indeed $f(Y_i^+)-f(Y_i^-)$ integrates $f'$ over that interval, and replacing the
integrand by its value at an endpoint costs at most $D_i^2\operatorname{Lip}(f')$.
For Kolmogorov tests the step in the Stein derivative also introduces an interval-crossing
term weighted by $D_i$. The exact unweighted crossing identity
\[
 \sum_i h_i v_i\Prb(Y_i^-\le x<Y_i^+)=\Cov(Y,\ind_{\{Y>x\}})\le\sqrt V/2
\]
does not supply this length-weighted estimate.

\paragraph{An explicit sufficient criterion for Kolmogorov distance.}
Set $\alpha_i=h_iv_i$, and define the two functions
\[
 B(t)=\sum_i\alpha_i\Cov(D_i,\ind_{\{\widehat Y_i\le t\}}),\qquad
 J(t)=\sum_i\alpha_i\E[D_i\ind_{\{Y_i^-\le t<Y_i^+\}}].
\]
Then
\begin{equation}\label{e:criterion}
 \dK((Y-\E Y)/\sqrt V,\Nlaw)\le \frac{2}{V}
                  (\|B\|_\infty+\|J\|_\infty),\qquad
 \int_{\mathbb R}J(t)\,dt\le3WV .
\end{equation}
To prove this, weight the $i$th coupling by $\alpha_iD_i/V$ and sum over $i$;
the resulting measure has total mass one. Sample the anchored endpoint and a uniform
point on the interval in this measure. The endpoint has cdf $F_Y+B/V$. The cdf
$F_*$ of the uniform point differs from the endpoint cdf at $t$ by at most $J(t)/V$.
By \eqref{e:identity} this point has the zero-bias law translated back to mean
coordinate $\E Y$: writing it $Y^*$, one has $\Cov(Y,f(Y))=V\E f'(Y^*)$.
For the standardized normal Stein solution the total variation of $f'_x$ is at
most two. Indeed $u f_x(u)$ increases continuously from $-(1-\Phi(x))$ to $\Phi(x)$:
on $u\le x$ it equals $(1-\Phi(x))[-1+\int_0^\infty t e^{ut-t^2/2}\,dt]$,
and on $u\ge x$ it equals $\Phi(x)[1-\int_0^\infty t e^{-ut-t^2/2}\,dt]$.
The additional step in $f'_x$ has variation one. Integration against the cdf
difference thus gives $\dK((Y-\E Y)/\sqrt V,\Nlaw)\le2\|F_Y-F_*\|_\infty$.
The integral identity $\int J=\sum_i\alpha_i\E D_i^2$ and Theorem~\ref{t:cost}
finish \eqref{e:criterion}.
In particular the two concrete bounds $\|B\|_\infty,\|J\|_\infty\le CW\sqrt V$
would suffice. Splitting at a length $l>0$ gives the weaker estimate
$J(t)\le l\sqrt V/2+3WV/l$ from the two crossing estimates above; optimizing gives
$\|J\|_\infty/V\le\sqrt{6W/\sqrt V}$.

\begin{prop}[A fluctuation obstruction for the direct switch]\label{p:obstruct}
Write $M(\eta)=\sum_i h_i v_i\E[D_i\mid\widehat\eta_i=\eta]$ under the
original clock law. Then $\E M=V$, but even for unit-weight pairs
$\Var(\E[M\mid Y])/(W^2 V)$ can grow at least as a constant times $\sqrt r$.
For real nonnegative weights the universal bound
$\E|\E[M\mid Y]-V|\le CW\sqrt V$ also fails for this switch.
\end{prop}
\begin{proof}
Use $m=r^4$ threshold-one leaves, one hub of threshold $r$, and pair clock
mean $a=r/m$ at every leaf. The occupied leaf count has law
$\operatorname{Bin}(m,1-e^{-a})$, and the additional hub is bounded, so $V\sim r$.
Let $E$ be the event of $r-1$ arrivals to distinct leaves. The probability of any
repeat is at most $ma^2/2=1/(2r^2)$. With no repeat, $Y=r-1$ occurs precisely on
$E$, hence $\Prb(Y=r-1)\sim(2\pi r)^{-1/2}$ and $\Prb(E\mid Y=r-1)\to1$.
On $E$ each absent leaf is the lower endpoint of its conditional switch. The
upper endpoint necessarily activates that leaf and the hub, so $D_i=2$.
Positivity gives $M\ge2(m-r+1)(1-e^{-a})e^{-a}=(2-o(1))r$ on $E$.
This atom gives $\Var(\E[M\mid Y])\ge(1-o(1))r^{3/2}/\sqrt{2\pi}$.

For the last assertion take $k$ independent copies of the star at fixed $r$, scaling
all weights in copy $b$ by $c_b=1+t^b$ for a fixed transcendental $0<t<1$.
The weights $c_b$ are linearly independent over the rationals, so the real count
$\sum_b c_bY_b$ determines all the integer counts $Y_b$. The new conditional
mass minus its mean is $\sum_b c_b^2(\E[M_b\mid Y_b]-V)$ with independent terms.
The new variance is $V\sum_b c_b^2$ and the batch bound is at most four. At fixed $r$
these terms are bounded uniformly in $b$; as $k\to\infty$, the CLT and bounded
normalized second moments give mean absolute deviation asymptotic to $\sqrt{2/\pi}$
times the standard deviation of this sum. The ratio to the new $W\sqrt{\Var(Y_{\rm total})}$
is at least a constant times $r^{1/4}$ for large $r$ and then large $k$.
Theorem~\ref{t:block} applies throughout this construction.
\end{proof}
Thus taking an absolute value on the conditional-mass fluctuation loses information
even in a class with the optimal normal approximation. For unrestricted pair graphs,
the unresolved estimates concern the \emph{signed} covariance in \eqref{e:gap}
tested against the normal Stein solution and its length-weighted crossing error.

\begin{thebibliography}{9}
\bibitem{previous} \emph{An optimal variance rate for Poissonized higher coverage},
\href{https://github.com/long-mathematics/bounded-batch-discovery/blob/main/papers/higher_coverage.pdf}{companion research note}, November 27, 2026. See also the
\href{https://github.com/long-mathematics/bounded-batch-discovery}{bounded-batch-discovery}
manuscripts for the Poissonized model, first-discovery bound, and threshold envelope.
\bibitem{BC} S. G. Bobkov and G. P. Chistyakov, \emph{On concentration functions of random
variables}, J. Theor. Probab. 28 (2015), 976--988, equations (1.2)--(1.3).
\href{https://www-users.cse.umn.edu/~bobko001/papers/2015_JOTP_BC_Conc.functions.authors.03.07.2013.pdf}{Article}.
\bibitem{MR} A. L. Miroshnikov and B. A. Rogozin, \emph{Inequalities for concentration
functions}, Theory Probab. Appl. 25 (1980), 176--180.
\end{thebibliography}
\end{document}
